\section{Definición operacional de la AGI: cinco capacidades clave}

\subsection{Por qué se necesita una nueva definición}

El académico Zhang Bo (张钹) señala que las definiciones actuales de la AGI (como la de Musk, «completar más del 70\% de las tareas humanas») son \textbf{completamente inejecutables e inverificables}, y necesariamente dan lugar a controversia.

\begin{warningbox}{El daño de las definiciones ambiguas}
Tomemos como ejemplo el reconocimiento de patrones: ¿basta con que la tasa de reconocimiento supere a la de un humano para considerar que se «alcanzó el nivel humano»? Algunos dicen que sí, otros dicen que, visto desde la robustez, todavía hay una distancia considerable. Una definición ambigua hace que la discusión sobre la AGI caiga en una polémica interminable.
\end{warningbox}

\subsection{Cinco capacidades clave}

El académico Zhang Bo propone una definición de la AGI \textbf{ejecutable y verificable}, es decir, que debe alcanzar las siguientes cinco capacidades clave (presten especial atención al \textbf{adjetivo calificativo} de cada una):

\begin{importantbox}{Las cinco capacidades clave de la AGI}
\begin{enumerate}
\item Comprensión multimodal \textbf{consistente en el espacio y el tiempo}: la clave está en la «consistencia espaciotemporal»; los distintos modos tienen ritmos temporales distintos (el video avanza cuadro por cuadro, el texto puede resumir milenios en una frase), lo que hace extremadamente difícil alinearlos
\item Aprendizaje y adaptación \textbf{en línea}: se distingue del aprendizaje fuera de línea; el reto central es la \textbf{controlabilidad}, es decir, si el proceso de aprendizaje por refuerzo puede converger hacia el objetivo
\item Razonamiento \textbf{verificable} y planeación \textbf{de largo plazo}: el resultado del razonamiento debe poder verificarse; la planeación debe tener carácter de largo plazo
\item Reflexión y metacognición \textbf{verificables}: la reflexión actual es de naturaleza «sensible», carece de una señal rastreable y verificable
\item Generalización \textbf{entre tareas}: no basta con la generalización entre dominios (en la que los modelos grandes ya se desempeñan relativamente bien), sino que también debe ser una generalización entre tareas fuera de distribución, con estructuras distintas y en escenarios de cola larga
\end{enumerate}
\end{importantbox}

\begin{knowledgebox}{Por qué importan los adjetivos}
El académico Zhang Bo insiste repetidamente en que «el adjetivo es muy importante». Por ejemplo: muchos equipos trabajan en «comprensión multimodal», pero al agregar «consistente en el espacio y el tiempo» la dificultad aumenta drásticamente; muchos modelos tienen «razonamiento», pero el razonamiento «verificable» es una exigencia completamente distinta. Estos calificativos delimitan la verdadera frontera técnica.
\end{knowledgebox}

